\hypertarget{system-analysis-and-requirements}{%
\chapter{System Analysis and
Requirements}\label{system-analysis-and-requirements}}

\begin{draftnote}

\textbf{Draft status, 04 November 2026.} This chapter is the narrative
summary; the full specification is Appendix A
(\texttt{docs/report/SRS.md}), which carries every requirement with a
stable identifier. Requirement numbers used here refer to that document
--- do not renumber them, because Chapter 6's verification refers to the
same identifiers.

\end{draftnote}

\begin{center}\rule{0.5\linewidth}{0.5pt}\end{center}

\hypertarget{the-problem-stated-as-a-requirement}{%
\section{The problem, stated as a
requirement}\label{the-problem-stated-as-a-requirement}}

A member of faculty marking fifty or five hundred programming
submissions needs to know \textbf{where to spend attention}. They do not
need a verdict, and a tool is not entitled to give one (Chapter 2 \S{}2.2).
What they need is an ordering, and a reason for the ordering they can
read, defend to a student, and be held to.

That reframing drives nearly every requirement that follows. The
system's output is a \textbf{ranking with attached evidence}, not a
percentage with a threshold.

\hypertarget{existing-systems-and-where-they-stop}{%
\section{Existing systems, and where they
stop}\label{existing-systems-and-where-they-stop}}

\begin{longtable}[]{@{}
  >{\raggedright\arraybackslash}p{(\columnwidth - 6\tabcolsep) * \real{0.2500}}
  >{\raggedright\arraybackslash}p{(\columnwidth - 6\tabcolsep) * \real{0.2500}}
  >{\raggedright\arraybackslash}p{(\columnwidth - 6\tabcolsep) * \real{0.2500}}
  >{\raggedright\arraybackslash}p{(\columnwidth - 6\tabcolsep) * \real{0.2500}}@{}}
\toprule\noalign{}
\begin{minipage}[b]{\linewidth}\raggedright
system
\end{minipage} & \begin{minipage}[b]{\linewidth}\raggedright
approach
\end{minipage} & \begin{minipage}[b]{\linewidth}\raggedright
what it finds
\end{minipage} & \begin{minipage}[b]{\linewidth}\raggedright
what it cannot find
\end{minipage} \\
\midrule\noalign{}
\endhead
\bottomrule\noalign{}
\endlastfoot
MOSS & winnowing fingerprints over k-grams & submissions resembling each
other & anything with no second copy in the corpus \\
JPlag & language-aware tokenising, greedy string tiling & the same & the
same \\
Turnitin and similar & text similarity against a web and paper corpus &
copied prose, some copied code & code transformed by renaming or
reordering \\
\end{longtable}

All of the above compare submissions \textbf{against each other}. That
is sufficient for collusion, which is the most common case, and
structurally incapable of detecting work outsourced to a third party or
generated by a model, because neither leaves a second copy in the
cohort.

\textbf{The proposed system differs in asking a second question}: not
\emph{does this resemble another submission?} but \emph{does this
resemble its own author's previous work?} Chapter 7 reports the extent
to which the implemented answer to that second question succeeds, and it
is more limited than the design intended.

\hypertarget{user-classes}{%
\section{User classes}\label{user-classes}}

\begin{longtable}[]{@{}
  >{\raggedright\arraybackslash}p{(\columnwidth - 2\tabcolsep) * \real{0.5000}}
  >{\raggedright\arraybackslash}p{(\columnwidth - 2\tabcolsep) * \real{0.5000}}@{}}
\toprule\noalign{}
\begin{minipage}[b]{\linewidth}\raggedright
class
\end{minipage} & \begin{minipage}[b]{\linewidth}\raggedright
may
\end{minipage} \\
\midrule\noalign{}
\endhead
\bottomrule\noalign{}
\endlastfoot
\textbf{Student} & submit work; see their own status and any recorded
decision; answer a flag; nominate past work as an anchor \\
\textbf{Faculty} & create courses and assignments; enrol students; see
the ranked queue and full evidence; record decisions; accept
nominations; re-run detection \\
\textbf{Admin} & all of the above, plus create accounts with roles and
edit detection configuration \\
\end{longtable}

\textbf{A student never sees a score.} This is a requirement (FR-30),
not an oversight: a ranking a student can read is a ranking a student
can tune against. The endpoint backing the student view returns no score
at all, so the property is enforced by the API rather than by the
interface choosing not to display it.

\hypertarget{functional-requirements-in-summary}{%
\section{Functional requirements, in
summary}\label{functional-requirements-in-summary}}

Thirty-two functional requirements, grouped. The full statements are in
Appendix A.

\begin{itemize}
\tightlist
\item
  \textbf{Authentication and accounts (FR-1 to FR-6).} Public
  self-registration creates a student account and no other; role
  assignment is an admin operation.
\item
  \textbf{Courses, enrolment, configuration (FR-7 to FR-10).} Detection
  configuration is per course, with per-language overrides --- because
  the evidence in Chapter 7 \S{}7.4 shows the same feature behaving
  oppositely in Python and Java, a single configuration for both cannot
  be correct.
\item
  \textbf{Assignments (FR-11, FR-12).} An assignment carries its
  language and its provenance, invigilated or take-home, because
  provenance determines whether the work may later serve as a style
  anchor.
\item
  \textbf{Submission (FR-13 to FR-16).} One source file,
  extension-checked and size-limited, stored with a content hash, with
  detection enqueued and debounced rather than run inline.
\item
  \textbf{Detection (FR-17 to FR-22).} Both layers computed where
  possible, composed with renormalisation, with the applied weights, the
  threshold and both layers' \textbf{evidence} stored on the result.
\item
  \textbf{Baseline integrity (FR-23 to FR-28).} The five mitigations
  described in Chapter 4 \S{}4.6.
\item
  \textbf{Review and the student view (FR-29 to FR-31).}
\item
  \textbf{Audit (FR-32).} Every state change affecting a student's
  record, held in a log separate from the documents it describes.
\end{itemize}

\hypertarget{non-functional-requirements}{%
\section{Non-functional
requirements}\label{non-functional-requirements}}

Ten, of which three are worth naming here.

\textbf{Bounded comparison (NFR-3).} The detector accepts at most 50
candidates and compares at most 10 in full, selected by a cheap
prefilter. This bounds the cost of one analysis irrespective of class
size --- and it is also a constraint on \emph{measurement}: a study that
compares exhaustively is measuring a system that was never deployed.
That error was made and corrected (Chapter 7 \S{}7.2).

\textbf{Explainability (NFR-8).} Every feature contributing to a
behavioural score must be explainable to a student contesting a flag.
This is why four features are scored rather than the ten extracted, and
it is a genuine constraint on accuracy: a feature a reviewer cannot
explain is not evidence, whatever it contributes to a score.

\textbf{Stated limitations (NFR-9, NFR-10).} Two requirements in the
specification record what the system does \emph{not} establish --- Layer
2's efficacy against partial outsourcing, and Layer 1's precision and
recall on genuine plagiarism. Recording a limitation as a numbered
requirement is unusual, and it is deliberate: it makes the gap traceable
rather than something a reader must infer from an absence.

\hypertarget{feasibility}{%
\section{Feasibility}\label{feasibility}}

\textbf{Technical.} The parsing and tree-comparison libraries required
are mature and freely available; the architecture is conventional. The
one genuine technical risk was the cost of tree edit distance, and it is
addressed by the prefilter.

\textbf{Operational.} The system requires something a course may not
have: \textbf{attributable past work for each student.} Layer 2 cannot
function without at least three trustworthy anchors per student per
language, and where those do not exist the layer declines to score
rather than guessing. A course with no invigilated work and no
nomination process gets a one-layer system.

\textbf{Ethical.} The system ranks students for suspicion, which carries
a direct risk of harm to an honest student through a wrong baseline.
That risk is the reason for the five mitigations, for the requirement
that a student may answer a flag, for the separate audit log, and for
the decision that no output of this system constitutes a finding of
misconduct.

\textbf{Scope.} Layer 3 was excluded on requirements grounds rather than
schedule grounds, and the reasoning is in Chapter 5 \S{}5.4.

\hypertarget{verification-approach}{%
\section{Verification approach}\label{verification-approach}}

Requirements are verified by the 278-test suite described in Chapter 6,
with identifiers traceable from Appendix A to test cases. NFR-9 and
NFR-10 are verified \textbf{as limitations rather than as capabilities},
by the committed measurement scripts reported in Chapter 7 --- which is
the only honest way to verify a requirement that records what the system
cannot do.
